\section{Discussion and Implications}
\subsection{Theoretical Implications}

Each implication below depends on the propositions receiving support.

For ethical leadership and social learning research, the model makes an assumption visible. The case for role modeling rests on followers being able to see their models and the consequences of their models' conduct~\citep{brown2005ethical,brown2006ethical}. Observational access turns that assumption into a variable, and leader ethical narration identifies a path by which ethical leadership can work when observation is limited. This is consistent with the view that vicarious learning can take place through conversation as well as through watching~\citep{myers2018coactive}, and it suggests that the communicative side of ethical leadership may matter more, relative to silent example, when teams are distributed.

For socialization research, the model suggests that physical arrangements belong in the analysis of socialization tactics. Van Maanen and Schein's contrast between serial socialization, in which newcomers learn from predecessors, and disjunctive socialization, in which they have none to learn from~\citep{vanmaanen1979toward}, was framed in terms of whether predecessors exist. Distributed entry may produce a hidden form of disjunctive socialization in which predecessors exist but are seldom seen. If so, an organization could believe it socializes newcomers serially while newcomers experience something closer to the opposite.

For imprinting research, the model proposes that ethical norms can be imprinted in the same way that approaches to work and professional orientations have been shown to be~\citep{tilcsik2014imprint,azoulay2017social}. It also specifies an imprint defined partly by absence: what a newcomer did not observe at entry may shape later conduct as much as what the newcomer did observe.

For generational research, the model offers a way to keep the idea of a formative period without relying on birth-year categories. Critics have argued that birth cohorts explain little about work attitudes~\citep{costanza2015generationally,rudolph2017considering}. Defining the relevant cohort by shared conditions of entry~\citep{joshi2010unpacking} yields claims that can be tested directly, because entry conditions can be measured for newcomers of any age.

For behavioral ethics, the two-sided effect in Proposition 3 implies that the ethical consequences of distributed work should differ across units. Research that estimates an average effect of remote work on misconduct could find little or nothing even if the model is correct, because losses in ethical units and partial protection in units with normalized misconduct would offset each other.

\subsection{Practical Implications}

The suggestions that follow are hypotheses for practice derived from the propositions, not recommendations backed by evidence.

Supervisors of distributed newcomers may need to say out loud what co-located newcomers would have seen. Explaining why a borderline request was refused, or what followed when a rule was ignored, is the behavior the moral manager concept describes~\citep{trevino2000moral}, and in our model it is the main available substitute for observation. Making the consequences of conduct visible, which ethical culture research treats as one dimension of an ethical organization~\citep{kaptein2008developing}, serves the same purpose, although it has to be balanced against the confidentiality owed to the people involved.

Organizations that rely on codes and online training to onboard remote newcomers should not assume that these convey the practiced norms, since those formal elements are distinct from the informal ones that carry enacted norms~\citep{tenbrunsel2003building}. Where in-person time is limited, scheduling part of it early in the entry period, and around work in which ethical judgment is exercised, is one way to raise observational access when it matters most. That timing recommendation is our inference from imprinting theory, not an established finding. Structured, shared entry experiences of the kind associated with institutionalized socialization tactics~\citep{ashforth1996socialization} may also help. Because newcomers obtain normative information largely from peers~\citep{morrison1993newcomer}, pairing distant newcomers with well-chosen peers is another option.

Newcomers themselves are not passive in this process. Remote newcomers have been observed seeking out insiders to make up for the contact they lacked~\citep{paik2026membership}. The model suggests that the most useful questions for a distant newcomer are about cases rather than rules: how a particular difficult request was handled, and what happened to the people involved. Such questions still carry image costs~\citep{ashford2003reflections}, which is one reason the burden should not rest on newcomers alone.

Finally, the two-sided effect has an uncomfortable implication. In units where misconduct has become normal, distance may partly protect newcomers, but it does nothing to change the unit. Such units require changes to the enacted norms themselves~\citep{ashforth2003normalization}.

\subsection{A Testing Agenda}

The propositions can be tested with established designs, but new measures are needed first. Observational access, enacted-norm clarity, and leader ethical narration have no validated measures. Enacted-norm clarity could be assessed by asking newcomers how conduct in specific gray-area situations is actually handled in their unit and comparing their answers with those of experienced insiders. Observational access could be measured by asking how often newcomers have witnessed colleagues or supervisors deal with ethically charged situations and learned what followed. Existing instruments can cover the other constructs: ethical leadership~\citep{brown2005ethical}, ethical culture including its transparency dimension~\citep{kaptein2008developing}, coworker antisocial behavior as an indicator of normalized misconduct~\citep{robinson1998monkey}, newcomer information seeking~\citep{morrison1993newcomer}, and professional isolation~\citep{golden2008impact}.

The most informative field design would compare distributed and co-located entrants within the same organizations. Newcomers are not randomly assigned to distributed entry, and remote roles may differ from on-site ones in function and autonomy, so such comparisons need to address selection. Changes in onboarding policy, such as a move from remote to on-site onboarding at a given date, create entry cohorts that differ in their conditions of entry while sharing an organization, a logic similar to that of studies using organization-wide changes in work arrangements~\citep{yang2022effects}. Measuring enacted-norm clarity during entry and conduct afterwards would test Propositions 1 to 3. Following both groups after their working arrangements converge would test Proposition 4. Including older entrants, such as career changers, in both conditions is what makes Proposition 5 testable.

Because newcomers are nested in work units, and Proposition 3 involves the norms of the unit, the natural analytic framework is multilevel, with the ethical quality of enacted norms measured at the unit level and its interaction with enacted-norm clarity or distributed entry tested across levels. The entry period also needs an explicit definition in each study, for example the first several months of tenure, and enacted-norm clarity should be measured more than once within it, since the model concerns how quickly and how completely newcomers learn the enacted norms.

Scenario experiments can complement field work. Participants could be shown the same work situation with or without the opportunity to observe how a colleague's shortcut was handled, and with or without a supervisor's explanation, and then asked what the unit's norm is and what they would do. Such designs support causal inference about Propositions 1 and 2, but they measure inferences and intentions rather than conduct, so field evidence would still be needed. Research on normalized misconduct also raises its own ethical questions, and studies of Proposition 3 would need strong protections for participants.

\subsection{Boundary Conditions and Ways the Model Could Fail}

Several conditions should weaken the proposed relationships. Structured socialization programs that give distributed newcomers regular, shared contact with insiders may preserve observational access despite distance~\citep{ashforth1996socialization}. Newcomers with substantial prior work experience bring ethical schemas formed elsewhere, and research on career history suggests that such prior experience carries over~\citep{dokko2009unpacking,higgins2005career}; for them, the entry period in a new organization should matter less than it does for people entering a first professional job. Newcomers with strong personal ethical commitments may resist normalized misconduct even when they see it clearly, which would weaken Proposition 3 in units where misconduct is normal.

The details of distributed arrangements matter too. Teams that work largely through synchronous video and open channels may give newcomers more visibility than their physical separation implies, in line with evidence that visible digital communication lets observers learn about colleagues~\citep{leonardi2015ambient}. Hybrid arrangements also vary in whether newcomers and the people they learn from are in the office on the same days, which may matter more than the number of days either spends there. Occupational and national context may matter as well. Where professional bodies or regulators set detailed rules of conduct, formal norms may leave less room for the unit's enacted norms, and the model's effects should be smaller.

The model could also be wrong, and its propositions are designed so that this can be detected. The main rival is a monitoring account, in which distance reduces supervision and so increases the opportunity for misconduct in every unit. That account predicts that distributed entrants behave less ethically regardless of their unit's norms. Our model predicts a crossover, with distributed entrants behaving less ethically than co-located entrants in ethical units but adopting less misconduct in units where it is normal. If the crossover does not appear, Proposition 3 fails and the monitoring account gains support. If differences between distributed and co-located entrants disappear once their working arrangements converge, Proposition 4 fails, which would suggest that ethical learning lost at entry can be recovered later. If members of Generation Z differ from older newcomers who entered under the same conditions, Proposition 5 fails, and some cohort effect beyond entry conditions would need explaining.

\subsection{Limitations}

The paper reports no data, and none of its propositions has been tested. Three of its constructs are new and unmeasured. Distributed entry covers a wide range of arrangements, and the model treats it as a single dimension. We do not estimate how many members of Generation Z experienced distributed entry, and that share probably differs across occupations and countries. In our reading, research on remote onboarding is still thin, and most of the literature the model draws on comes from North American and European organizations. Measuring ethical conduct is difficult in any study, and self-reports of it are vulnerable to socially desirable responding.
